% ===================================================================
% Front matter: identity, rules status, legends, contents (2 pages)
% ===================================================================
\thispagestyle{plain}
\pdfbookmark[0]{Front: identity and rules status}{front}
\begin{center}
{\Huge\bfseries Thermodynamics \textendash\ Division C}\\[3pt]
{\Large Invitational Reference Binder (unofficial study reference)}\\[4pt]
{\normalsize Built \builddate{} for a tournament on \textbf{Thursday, 22 October 2026} (``tomorrow'' resolved in America/Chicago; adjust if your time zone differs)}
\end{center}
\vspace{2pt}
{\small
\begin{tabularx}{\textwidth}{@{}lXlX@{}}
Team / school: & \hrulefill & Tournament: & \hrulefill\\[4pt]
Partner A: & \hrulefill & Partner B: & \hrulefill\\
\end{tabularx}
Suggested split (adaptable): A --- Tabs 3--4 (gases, laws, cycles); B --- Tabs 2, 5, 7 (heat, transfer, data); share Tabs 1 and 6.}
\vspace{2pt}

\begin{rulebox}\label{sec:status}\textbf{RULES STATUS --- READ BEFORE THE EVENT.} The official 2027 Thermodynamics rules, official clarifications, and this tournament's instructions were \textbf{not inspected}: none were attached, and the official site was blocked from the build environment. Therefore the following are \textbf{UNVERIFIED}:
\begin{enumerate}
\item whether this binder may be used at all, and in what form (page or sheet limit, binding, protectors, tabs, two-sided printing);
\item which calculator class is allowed;
\item the written-test scope or topic list (topic labels below are editorial, not official);
\item device construction, procedure, timing, and scoring (\pg{sec:device});
\item any modifications announced by the tournament host.
\end{enumerate}
Snippet-level information only (not verified): the official event description reportedly reads ``construct a device prior to the tournament designed to \emph{collect heat} and complete a written test on thermodynamic concepts,''\src{SO1} while a community wiki describes a \emph{heat-retention} beaker task.\src{WIKI} The conflict is unresolved. \textbf{An unverified permission is not a scientific error:} the physics in this binder was checked independently (\pg{sec:sources}); its \emph{eligibility} for use was not.\end{rulebox}

\begin{multicols}{2}
\begin{reftab}{T:alignment}{Alignment summary (full matrix in \texttt{rules-alignment.md}).}
\small
\begin{tabularx}{\linewidth}{@{}L{0.33\linewidth}L{0.24\linewidth}Y@{}}
\toprule
Item & Status & Location\\
\midrule
Written-test scope & unverified & whole binder; editorial labels\\
Reference / calculator permissions & unverified & this page\\
Device participation and limits & unverified; sources conflict & \pg{sec:device}\\
Event in 2027 Div.\ C slate & snippet only & \pg{sec:sources}\\
\bottomrule
\end{tabularx}\end{reftab}

\textbf{Pre-event checklist (most important first)}
\begin{enumerate}
\item[$\square$] Confirm in the 2027 Rules Manual (and host instructions) that a binder is allowed, and its size and format limits. Remove pages if required.
\item[$\square$] Confirm the calculator class.
\item[$\square$] Confirm device specs, impound time, energy source, and what is scored before touching the device (\pg{sec:device}).
\item[$\square$] Check for official clarifications/FAQs and host modifications.
\end{enumerate}

\columnbreak
\textbf{Content labels (editorial, not official):}\\
\CORE{} prioritized foundation; \EXT{} useful added depth, no claim it is tested; \RD{} construction, scoring, permissions, or host-specific.

\textbf{Box types:}
\begin{formulabox}canonical equations with IDs.\end{formulabox}
\begin{keybox}model choice and decision steps.\end{keybox}
\begin{demobox}short, fully solved, numerically checked.\end{demobox}
\begin{trapbox}common mistakes and validity limits.\end{trapbox}

\textbf{Source tags} (superscript, clickable): e.g.\ \srcin{IAPWS}, \srcin{CODATA}, \srcin{DER} (derived/checked here), \srcin{EK} (editor knowledge, unverified this build). Register: \pg{sec:sources}.

\textbf{Conventions:} $\Delta U=Q-\Wby$ (\pg{T:signs}); kelvin in ratios and powers; pressures absolute; cycle quantities are magnitudes.
\end{multicols}

\clearpage
\pdfbookmark[0]{Contents}{contents}
{\small
\renewcommand{\contentsname}{Contents}
\makeatletter
\renewcommand\tableofcontents{%
  \section*{\contentsname}\vspace{-4pt}%
  \begin{multicols}{2}\@starttoc{toc}\end{multicols}}
\makeatother
\tableofcontents}
\vspace{4pt}
\textbf{How to use:} (1) Find the cue in the \textbf{question router} (\pg{sec:router}) or the \textbf{index} (\pg{sec:index}). (2) For a calculation, use the \textbf{formula selector} (\pg{sec:selector}) to confirm the equation's validity and hazards, then jump to its canonical page. (3) Data tables: \tref{T:datalocator}, \pg{T:datalocator}. Equation and table numbers start with the tab number (Eq.\ (3.20) is in Tab~3).
